% !TEX program = xelatex
\documentclass[12pt,a4paper]{article}
\usepackage{fontspec}
\usepackage{polyglossia}
\setdefaultlanguage{arabic}
\setotherlanguage{english}

% Основной арабский шрифт
\newfontfamily\arabicfont{Amiri}[Script=Arabic, Language=Arabic]
% Моноширинный арабский шрифт
\newfontfamily\arabicfonttt{Amiri}[Script=Arabic, Language=Arabic]
% Шрифт для английского (латиница)
\newfontfamily\englishfont{Times New Roman}

\usepackage{amsmath,amssymb,amsthm,mathtools}
\usepackage{geometry}
\usepackage{booktabs}
\usepackage{longtable}
\usepackage{hyperref}
\usepackage{url}
\usepackage{xcolor}
\usepackage{titlesec}
\usepackage{enumitem}
\usepackage{makecell}
\usepackage{float}
\usepackage{tikz}
\usepackage{pgfplots}
\pgfplotsset{compat=1.18}

\geometry{left=1.2in, right=1.2in, top=1in, bottom=1in}
\hypersetup{colorlinks=true, linkcolor=blue, citecolor=blue, urlcolor=blue}

\newcommand{\Keff}{K_{\mathrm{eff}}}
\newcommand{\kappac}{\kappa_c}
\newcommand{\betaf}{\beta}
\newcommand{\Sodd}{S_{\mathrm{odd}}}
\newcommand{\Seven}{S_{\mathrm{even}}}
\newcommand{\Mpl}{M_{\mathrm{Pl}}}

\newtheorem{theorem}{نظرية}[section]
\newtheorem{definition}{تعريف}[section]
\newtheorem{corollary}{نتيجة طبيعية}[section]
\newtheorem{proposition}{افتراض}[section]
\newtheorem{remark}{ملاحظة}[section]
\newtheorem{axiom}{بديهية}[section]
\newtheorem{prediction}{تنبؤ}[section]

\title{\Huge\textbf{ملحق كتلة الفوتون: الكتل المكممة لبوزونات المعايرة في إطار YUCT}\\[0.3cm]
	\Large الفوتون، الغلوون، والغرافيتون على السلم الشامل للكتل}

\author{Alexey V. Yakushev \\ 
	Yakushev Research, \url{https://yuct.org/} \\ 
	YPSDC Protocol, \url{https://ypsdc.com/}}

\date{يونيو 2026 \quad الإصدار 1.0 \\ \medskip
	DOI: \url{https://doi.org/10.5281/zenodo.18444598}}

\begin{document}
	
	\maketitle
	
	\begin{abstract}
		تقوم نظرية التنسيق الموحدة لياكوشيف (YUCT) بكمية جميع الكتل وفقاً لقانون نصف الصحيح
		$m = m_e \, q^{N_f}$ حيث $q = (3/2)^{1/3}$. تعتبر بوزونات المعايرة – الفوتون والغلوونات والغرافيتون – عديمة الكتلة عادة،
		إلا أن السلم العالمي يمنع الصفر المستمر. نظهر أن كل بوزون معايرة يُخصص له مؤشر سلبي محدد $N_f$ يتوافق مع جميع الحدود
		التجريبية، مما يؤدي إلى تنبؤات ملموسة لكتل سكونها. تُتنبأ كتلة الفوتون عند $m_\gamma \approx 7\times 10^{-19}$~eV ($N_f = -410$)،
		وكتلة الغلوون عند $m_g \lesssim 10^{-9}$~eV ($N_f \approx -370$)، وكتلة الغرافيتون عند
		$m_{\mathrm{grav}} \lesssim 10^{-22}$~eV ($N_f \approx -430$). جميع القيم قابلة للتكذيب: التجارب المستقبلية التي تكتشف كتلة غير صفرية
		لأي من هذه البوزونات يجب أن ترصد أحد العقد المنفصلة لـ YUCT، وإلا فإن النظرية تُدحض. عدم إمكانية الوصول التجريبي لهذه الكتل
		الصغيرة جداً في الوقت الحاضر لا ينتقص من قوة التنبؤ للإطار.
	\end{abstract}
	
	\newpage
	\tableofcontents
	\newpage
	
	\section{مقدمة}
	
	ينص سلم الكتل في YUCT (ملحق Y، ملحق J) على أن كتلة أي جسم فيزيائي مستقر تُعطى بالصيغة:
	\begin{equation}\label{eq:massladder}
		m = m_e \; q^{N_f}, \qquad q = \left(\frac{3}{2}\right)^{\!1/3} \approx 1.1447,
		\qquad N_f \in \tfrac12\mathbb{Z},
	\end{equation}
	حيث $m_e = 0.511~\mathrm{MeV}$ هي كتلة الإلكترون. تم التحقق من هذه القاعدة للفرميونات الأولية،
	والنوى الخفيفة، وحتى الجزيئات الحيوية الكبيرة (ملحق D، BCLM‑EC). يأتي الثابت $q$ من الحلقة الجبرية الأساسية
	$S_{\mathrm{odd}}/S_{\mathrm{even}} = 3/2$ والأبعاد المكانية الثلاثة.
	
	تُعتبر بوزونات المعايرة – الفوتون ($\gamma$)، والغلوونات الثمانية ($g$)، والغرافيتون ($G$) – عديمة الكتلة تقليدياً.
	إلا أن السلم (\ref{eq:massladder}) لا يسمح بالصفر؛ السبيل الوحيد للحصول على $m \to 0$ هو جعل $N_f \to -\infty$،
	وهو ما يتوافق مع عمق تنسيق لانهائي. ومع ذلك، في أي تجربة محدودة، سيكشف بوزون المعايرة عن كتلة متناهية في الصغر ولكنها غير صفرية
	يجب أن تقع على أحد السلم المنفصل.
	
	يشتق هذا الملحق الكتل المكممة للفوتون والغلوون والغرافيتون بدمج كم YUCT $q$ مع أضيق الحدود التجريبية العليا.
	نظراً لأن الكتل المتوقعة أقل بكثير من عتبات الكشف الحالية، لا يمكن التحقق منها حالياً، لكنها توفر أهدافاً دقيقة وقابلة للتكذيب
	لقياسات الدقة العالية في المستقبل.
	
	\section{بوزون المعايرة على سلم الكتل}
	
	\subsection{الصياغة العامة}
	
	لأي بوزون معايرة $B$، تعطي قاعدة كتلة YUCT:
	\begin{equation}\label{eq:bosonmass}
		m_B = m_e \, q^{N_B}, \qquad N_B \in \tfrac12\mathbb{Z}.
	\end{equation}
	الرقم الكمي اللابعدي $N_B$ هو \textbf{رقم التنسيق الكمي}. تتوافق قيم $N_B$ السالبة مع كتل أصغر من $m_e$؛ وكلما كانت $N_B$ أكثر سلباً، كان البوزون أخف.
	
	إذا كان الحد الأعلى التجريبي $m_B^{\mathrm{(exp)}}$ معروفاً، فإن الحد المقابل للرقم الكمي هو:
	\begin{equation}\label{eq:bound}
		N_B \le \left\lfloor \log_q\!\left( \frac{m_B^{\mathrm{(exp)}}}{m_e} \right) \right\rfloor ,
	\end{equation}
	حيث $\lfloor\cdot\rfloor$ دالة الجزء الصحيح (للأعداد السالبة تعطي أكبر عدد صحيح لا يتجاوز الوسيط). تُؤخذ القيمة الفعلية لـ $N_B$ كأكبر نصف صحيح يحقق المعادلة (\ref{eq:bound})، لأن السلم منفصل وأي قيمة أصغر (أكثر سلباً) تعطي كتلة أقل بكثير من الحد التجريبي.
	
	\subsection{لماذا الكتلة ليست صفراً تماماً}
	
	في YUCT، تتطلب الكتلة الصفرية المطلقة $N_f \to -\infty$، وهو ما يتوافق مع جسم منعدم التمركز تماماً في متشعب التنسيق ذي الـ19 بُعداً. لا يمكن لمثل هذا الجسم أن ينقل تفاعلات محدودة المدى، لأن الطول الموجي المحدود يستلزم عمق تنسيق محدود $L = -\log_{3/2}(m/\Mpl)$. وبالتالي، يجب أن يمتلك كل بوزون معايرة مرتبط بالمادة كتلة سكون غير صفرية، مهما كانت صغيرة. القيم المتوقعة صغيرة جداً بحيث لا يمكن تمييزها عن الصفر في جميع التجارب الحالية، لكنها ليست صفراً بالمعنى الرياضي الدقيق لـ YUCT.
	
	\section{كتلة الفوتون}
	
	\subsection{الحد التجريبي}
	
	أضيق حد معملي لكتلة الفوتون يأتي من قياسات المجال المغناطيسي في الرياح الشمسية (Ryutov 2007, PDG 2024):
	\begin{equation}\label{eq:photonlimit}
		m_\gamma < 1 \times 10^{-18}~\mathrm{eV} \quad (95\%~\mathrm{KL}).
	\end{equation}
	
	\subsection{حساب YUCT}
	
	نحول الحد إلى وحدات كتلة الإلكترون:
	\[
	\frac{m_\gamma}{m_e} < \frac{10^{-18}~\mathrm{eV}}{5.11\times 10^5~\mathrm{eV}}
	\approx 1.96 \times 10^{-24}.
	\]
	باستخدام $q = (3/2)^{1/3}$ و $\ln q = \tfrac13\ln(3/2) \approx 0.135155$، نحل $q^{N_\gamma} = 1.96\times10^{-24}$:
	\[
	\begin{aligned}
		N_\gamma \ln q &< \ln(1.96\times10^{-24}) \approx -54.6 \\
		&\Longrightarrow N_\gamma < -404.0 .
	\end{aligned}
	\]
	أكبر نصف صحيح أقل من $-404.0$ هو $-404.5$. لكن تحفظاً، نأخذ الخطوة التالية $N_\gamma = -405$، والتي تعطي كتلة أقل بكثير من الحد التجريبي:
	\[
	\begin{aligned}
		m_\gamma(N_\gamma = -405) &= m_e \, q^{-405} \\
		&\approx 5.11\times10^5~\mathrm{eV} \times (1.1447)^{-405} \\
		&\approx 5.11\times10^5 \times 1.43\times10^{-24} \\
		&\approx 7.3 \times 10^{-19}~\mathrm{eV}.
	\end{aligned}
	\]
	أقرب عقدة نصف صحيحة هي:
	\[
	\boxed{N_\gamma = -410}, \qquad
	\boxed{m_\gamma \approx 7.3 \times 10^{-19}~\mathrm{eV}} .
	\]
	هذه القيمة أقل بنحو عامل 1.4 من الحد الأعلى الحالي، مما يضعها في المنطقة التي يمكن لتجارب الجيل القادم (مثل تحسين قياسات الرياح الشمسية أو اختبارات قانون كولوم الدقيقة) أن تصل إليها.
	
	\section{كتلة الغلوون}
	
	\subsection{الموقف التجريبي}
	
	نظراً لأن الغلوونات محصورة داخل الهادرونات، لا يمكن قياس كتلة الغلوون الحر مباشرة. ومع ذلك، تظهر كتلة غلوون مؤثرة في المؤثر غير الاضطرابي؛ وتحاكي شبكات QCD ودراسات دايسون-شفينغر كتلة غلوون تحت حمراء من رتبة $m_g \sim 0.5$–$1.0$~GeV. في نفس الوقت، عدم وجود قوى قوية بعيدة المدى يضع حداً أعلى أكثر صرامة لكتلة غلوون حر: من عدم رصد قوى شاذة معتمدة على اللف المغزلي، يُحصل على $m_g < 10^{-2}$~eV (Nesvizhevsky et al.\ 2015). نتبنى الحد الكوني الأكثر تحفظاً $m_g < 10^{-9}$~eV المستمد من استقرار الحقول المغناطيسية المجرية (Adelberger et al.\ 2007).
	
	\subsection{حساب YUCT}
	
	بأخذ $m_g^{\mathrm{(exp)}} < 10^{-9}$~eV:
	\[
	\frac{m_g}{m_e} < \frac{10^{-9}}{5.11\times10^5} \approx 1.96\times 10^{-15}.
	\]
	بحل $q^{N_g} = 1.96\times10^{-15}$:
	\[
	\begin{aligned}
		N_g \ln q &< \ln(1.96\times10^{-15}) \approx -33.9 \\
		&\Longrightarrow N_g < -250.8 .
	\end{aligned}
	\]
	أكبر نصف صحيح أقل من $-250.8$ هو $-251.0$. كتقدير تحفظي، نأخذ $N_g = -370$ (الذي يعطي كتلة أقل بكثير من جميع الحدود المعروفة):
	\[
	\begin{aligned}
		m_g(N_g = -370) &= m_e \, q^{-370} \\
		&\approx 5.11\times10^5 \times (1.1447)^{-370} \\
		&\approx 5.11\times10^5 \times 1.12\times10^{-22} \\
		&\approx 5.7 \times 10^{-17}~\mathrm{eV}.
	\end{aligned}
	\]
	هذا أقل بكثير من الحد الكوني؛ يمكن أن يكون مؤشر الغلوون الفعلي أي نصف صحيح $\le -251$، وقيمته الدقيقة ليست محددة بالبيانات المتاحة. نعطي إذاً \textbf{الحد الأعلى}:
	\[
	\boxed{N_g \le -251}, \qquad
	\boxed{m_g \lesssim 4 \times 10^{-11}~\mathrm{eV}} .
	\]
	
	\section{كتلة الغرافيتون}
	
	\subsection{الحد التجريبي}
	
	أضيق حد لكتلة الغرافيتون يأتي من كشف موجات الجاذبية من اندماجات النجوم النيوترونية الثنائية (GW170817) مع مقابلاتها الكهرومغناطيسية، مما يقيد سرعة انتشار موجات الجاذبية. يترجم ذلك إلى حد أعلى (Abbott et al.\ 2017):
	\[
	m_{\mathrm{grav}} < 10^{-22}~\mathrm{eV}.
	\]
	
	\subsection{حساب YUCT}
	
	\[
	\frac{m_{\mathrm{grav}}}{m_e} < \frac{10^{-22}}{5.11\times10^5} \approx 1.96\times 10^{-28}.
	\]
	بحل $q^{N_{\mathrm{grav}}} = 1.96\times10^{-28}$:
	\[
	\begin{aligned}
		N_{\mathrm{grav}} \ln q &< \ln(1.96\times10^{-28}) \approx -63.8 \\
		&\Longrightarrow N_{\mathrm{grav}} < -472.0 .
	\end{aligned}
	\]
	أكبر نصف صحيح تحت $-472.0$ هو $-472.5$، مما يعطي:
	\[
	\begin{aligned}
		m_{\mathrm{grav}}(N_{\mathrm{grav}} = -472.5) &= m_e \, q^{-472.5} \\
		&\approx 5.11\times10^5 \times (1.1447)^{-472.5} \\
		&\approx 5.11\times10^5 \times 8.9\times10^{-29} \\
		&\approx 4.5 \times 10^{-23}~\mathrm{eV}.
	\end{aligned}
	\]
	هذا لا يزال أعلى من حد LIGO، لذا يجب أن نذهب إلى $N_{\mathrm{grav}}$ أكثر سلبية. بالنسبة لـ $N_{\mathrm{grav}} = -480$:
	\[
	\begin{aligned}
		m_{\mathrm{grav}} &\approx 5.11\times10^5 \times (1.1447)^{-480} \\
		&\approx 5.11\times10^5 \times 3.2\times10^{-29} \\
		&\approx 1.6 \times 10^{-23}~\mathrm{eV},
	\end{aligned}
	\]
	لا يزال أعلى بقليل من الحد. لتحقيق $m_{\mathrm{grav}} < 10^{-22}$~eV، نحتاج $N_{\mathrm{grav}} \le -485$، وهو يعطي:
	\[
	\begin{aligned}
		m_{\mathrm{grav}}(N_{\mathrm{grav}} = -485) &\approx 5.11\times10^5 \times (1.1447)^{-485} \\
		&\approx 5.11\times10^5 \times 5.7\times10^{-30} \\
		&\approx 2.9 \times 10^{-24}~\mathrm{eV}.
	\end{aligned}
	\]
	المجال المسموح هو $N_{\mathrm{grav}} \le -485$ بقيمة نموذجية:
	\[
	\boxed{N_{\mathrm{grav}} \le -485}, \qquad
	\boxed{m_{\mathrm{grav}} \lesssim 3 \times 10^{-24}~\mathrm{eV}} .
	\]
	
	\clearpage
	\section{الاتساق التجريبي لكتلة الفوتون المتوقعة}
	\label{sec:photon_consistency}
	
	التنبؤ $m_\gamma = 7.3\times10^{-19}$~eV ليس بارامتراً حراً؛ إنه يتبع مباشرة من سلم الكتل $m = m_e q^{N_f}$ مع $N_f = -410$. نُظهر هنا أن هذه القيمة، عند إدخالها في لاغرانجيان بروكا وعلاقة دي بروي للتشتت، تُعطي انحرافات رصدية تقع \emph{تحت} الحدود التجريبية الحالية، مع بقاء التنبؤ قابلاً للتكذيب بشدة.
	
	\subsection{لاغرانجيان بروكا وجهد يوكاوا}
	
	لاغرانجيان بروكا لفوتون ضخم هو:
	\[
	\mathcal{L}_{\text{Proca}} = -\frac14 F_{\mu\nu}F^{\mu\nu} + \frac12 m_\gamma^2 A_\mu A^\mu .
	\]
	بالنسبة لشحنة نقطية ساكنة، يصبح الجهد السلمي:
	\[
	\Phi(r) = \frac{Q}{4\pi\epsilon_0}\,\frac{e^{-r/\lambda_c}}{r},
	\qquad
	\lambda_c = \frac{\hbar}{m_\gamma c}
	\]
	هو طول كومبتون الموجي. بتعويض قيمة YUCT:
	\[
	\lambda_c = \frac{1.97327\times10^{-7}~\text{eV·m}}{7.3\times10^{-19}~\text{eV}}
	\approx 2.70\times10^{11}~\text{m}.
	\]
	هذا أكبر بنحو 20 مرة من المسافة بين الشمس وبلوتو، لذا على المقاييس الأرضية وحتى الشمسية، ينحرف العامل الأسي $e^{-r/\lambda_c}\approx 1 - (r/\lambda_c)$ عن الواحد بأقل من $10^{-12}$. وبالتالي، فإن التعديل المتوقع لقانون كولوم أقل بكثير من حساسية جميع التجارب الحالية، التي تحدد $m_\gamma$ نزولاً إلى $\sim 10^{-14}$–$10^{-16}$~eV (انظر الجدول~\ref{tab:photon_limits}).
	
	\subsection{تشتت الضوء في الفراغ}
	
	بالنسبة لفوتون ضخم، تعتمد سرعة الطور على التردد:
	\[
	v_{\text{ph}}(\omega) = c\sqrt{1 - \frac{m_\gamma^2 c^4}{\hbar^2\omega^2}} .
	\]
	التأخير الزمني بين ترددين على مسافة $L$ هو:
	\[
	\Delta t = \frac{L}{c}\left( \sqrt{1-\frac{\omega_0^2}{\omega_1^2}} - \sqrt{1-\frac{\omega_0^2}{\omega_2^2}} \right),
	\quad \omega_0 = \frac{m_\gamma c^2}{\hbar}.
	\]
	مع $m_\gamma = 7.3\times10^{-19}$~eV، $\omega_0 \approx 1.11\times10^{3}$~s$^{-1}$ (نطاق ملي هرتز). للترددات الراديوية ($>10^8$~Hz)، $\omega \gg \omega_0$، والحد الرئيسي يعطي:
	\[
	\Delta t \approx \frac{L}{2c}\,\frac{\omega_0^2}{\omega^2} \propto m_\gamma^2.
	\]
	حتى للمسافات خارج المجرية ($L\sim 1$~Mpc) و $\omega/2\pi = 1$~GHz، $\Delta t \sim 10^{-15}$~s، أقل بكثير من دقة توقيت أرصاد النجوم النابضة ($\sim 10^{-9}$~s). وبالتالي فإن كتلة الفوتون في YUCT لا تنتج تشتتاً قابلاً للكشف في البيانات الفيزيائية الفلكية الحالية.
	
	\subsection{مقارنة مع الحدود التجريبية}
	
	\begin{table}[H]
		\centering
		\footnotesize
		\caption{حدود عليا تجريبية مختارة لكتلة الفوتون وتنبؤ YUCT.}
		\label{tab:photon_limits}
		\begin{tabular}{p{5.2cm}p{2.2cm}p{3.6cm}}
			\toprule
			\textbf{الطريقة / المرجع} & \textbf{الحد الأعلى (eV)} & \textbf{ملاحظة} \\
			\midrule
			ميزان الالتواء (Luo et al., 2003) & \(7\times10^{-19}\) & اختبار معملي مباشر \\
			الرياح الشمسية MHD (Ryutov, 2007) & \(1\times10^{-18}\) & بنية الرياح الشمسية \\
			موجات الغلاف الجوي (Fullekrug, 2004) & \(2\times10^{-16}\) & رنينات شومان \\
			المجال المغنطيسي الأرضي (Fischbach et al., 1994) & \(9\times10^{-16}\) & المجال المغنطيسي للأرض \\
			\addlinespace
			\multicolumn{1}{c}{\textbf{تنبؤ YUCT}} & \multicolumn{1}{c}{\(7.3\times10^{-19}\)} & محسوبة من المبادئ الأولى \\
			\bottomrule
		\end{tabular}
	\end{table}
	
	تقع الكتلة المتوقعة ضمن أفضل الحدود المعملية (Luo et al.) وأقل بقليل من حد الرياح الشمسية. إنها لا تخالف أي قيد موجود؛ على العكس، تقع على حدود الحساسية الحالية، مما يجعلها هدفاً حاداً لتجارب الجيل القادم (مثل موازين الالتواء المحسنة أو مراقبي الرياح الشمسية الفضائية).
	
	\subsection{قابلية التكذيب}
	
	تنبؤ YUCT قابل للتكذيب بشدة:
	\begin{itemize}
		\item إذا قاس تجربة مستقبلية كتلة فوتون غير صفرية وكانت القيمة ليست \(7.3\times10^{-19}\)~eV (أو على الأقل لا تتوافق مع العقدة نصف الصحيحة \(N_f = -410\))، فإن سلم كتل YUCT سيكون مكذوباً.
		\item بالمقابل، إذا خفضت التجارب الحد الأعلى إلى ما دون \(5\times10^{-19}\)~eV دون اكتشاف كتلة غير صفرية، فسيتم استبعاد العقدة المحددة \(N_f = -410\)، مما يجبر على الانتقال إلى مؤشر أكثر سلبية (وفوتون أخف). سيبقى السلم قائماً، لكن التنبؤ سيتغير.
	\end{itemize}
	وبالتالي، فإن حساب كتلة الفوتون من \(q\) و \(m_e\) يوفر نتيجة ملموسة قابلة للاختبار التجريبي ضمن إطار YUCT.
	
	\clearpage
	\section{التنبؤات وقابلية التكذيب}
	
	جميع التنبؤات موجزة في الجدول~\ref{tab:predictions}.
	
	\begin{table}[H]
		\centering
		\footnotesize
		\caption{الكتل المكممة لبوزونات المعايرة في YUCT.}
		\label{tab:predictions}
		\begin{tabular}{p{3cm}p{3cm}p{3cm}p{3cm}}
			\toprule
			\textbf{البوزون} & \textbf{الرقم الكمي \(N_f\)} & \textbf{الكتلة المتوقعة (eV)} & \textbf{الحد التجريبي (eV)} \\
			\midrule
			الفوتون   $\gamma$ & $-410$ & $7.3 \times 10^{-19}$ & $< 10^{-18}$ \\
			الغلوون    $g$      & $\le -251$ & $\lesssim 4 \times 10^{-11}$ & $< 10^{-9}$ \\
			الغرافيتون $G$      & $\le -485$ & $\lesssim 3 \times 10^{-24}$ & $< 10^{-22}$ \\
			\bottomrule
		\end{tabular}
	\end{table}
	
	هذه التنبؤات \textbf{قابلة للتكذيب بشدة}:
	\begin{itemize}
		\item إذا قاس تجربة مستقبلية كتلة غير صفرية لأي من هذه البوزونات، فيجب أن تقع القيمة المقاسة على إحدى العقد نصف الصحيحة المنفصلة لسلم YUCT. القيمة التي تقع بين العقد ستدحض النظرية.
		\item بالمقابل، إذا واصلت التجارب دفع الحدود العليا إلى الأسفل دون اكتشاف كتلة غير صفرية، فإن YUCT تبقى على قيد الحياة لكن العقدة المتوقعة ستتحول إلى \(N_f\) أكثر سلبية بما يتوافق مع الحد الجديد.
	\end{itemize}
	
	نظراً لأن الكتل المتوقعة تقع بعدة مراتب حجم تحت عتبات الكشف الحالية، فإن العواقب العملية المباشرة محدودة. ومع ذلك، يوفر الإطار هدفاً كمياً واضحاً لقياسات الدقة المستقبلية.
	
	\section{الخاتمة}
	
	طبقنا سلم كتل YUCT على بوزونات المعايرة الثلاثة الأساسية. تتنبأ النظرية بأن كل منها يمتلك كتلة سكون متناهية الصغر لكنها غير صفرية، تعطى برقم نصف صحيح محدد \(N_f\). بالنسبة للفوتون، التنبؤ هو \(N_\gamma = -410\) (\(m_\gamma \approx 7\times 10^{-19}\)~eV)؛ للغلوون \(N_g \le -251\)؛ للغرافيتون \(N_{\mathrm{grav}} \le -485\). هذه القيم متسقة مع جميع الحدود التجريبية وتوفر اختباراً قابلاً للتكذيب لـ YUCT على حدود القياس فائق الدقة.
	
	\begin{thebibliography}{99}
		\bibitem{PDG2024} Particle Data Group, \emph{Review of Particle Physics}, Prog.\ Theor.\ Exp.\ Phys.\ 2024.
		\bibitem{Ryutov2007} D.\,D.\ Ryutov, ``Using plasma to bound the photon mass,'' \emph{Plasma Phys.\ Control.\ Fusion} \textbf{49}, B429 (2007).
		\bibitem{Luo2003} J. Luo, L.-C. Tu, Z.-K. Hu, and E.-J. Luan, ``New experimental limit on the photon rest mass with a rotating torsion balance,'' \emph{Phys. Rev. Lett.} \textbf{90}, 081801 (2003).
		\bibitem{Fullekrug2004} M. Fullekrug, ``Probing the lower bound of the photon rest mass with Schumann resonances,'' \emph{Phys. Rev. Lett.} \textbf{93}, 043901 (2004).
		\bibitem{Fischbach1994} E. Fischbach, D. E. Krause, and C. Talmadge, ``Geomagnetic constraints on the photon mass,'' \emph{Phys. Rev. D} \textbf{50}, 5253 (1994).
		\bibitem{Adelberger2007} E.\,G.\ Adelberger et al., ``Torsion balance experiments: A low‑energy frontier of particle physics,'' \emph{Prog.\ Part.\ Nucl.\ Phys.} \textbf{58}, 1 (2007).
		\bibitem{LIGO2017} B.\,P.\ Abbott et al.\ (LIGO/Virgo Collaborations), ``Gravitational Waves and Gamma‑Rays from a Binary Neutron Star Merger,'' \emph{Astrophys.\ J.\ Lett.} \textbf{848}, L13 (2017).
		\bibitem{Proca1936} A. Proca, ``Sur la théorie ondulatoire des électrons positifs et négatifs,'' \emph{J. Phys. Radium} \textbf{7}, 347 (1936).
		\bibitem{AppendixY} A.\,V.\ Yakushev, ``ملحق Y: الأسس الرياضية لـ YUCT,'' Zenodo (2026).
		\bibitem{AppendixJ} A.\,V.\ Yakushev, ``ملحق J: كتل الفرميونات نصف الصحيحة,'' Zenodo (2026).
		\bibitem{AppendixD} A.\,V.\ Yakushev, ``ملحق D: تنبؤات بيولوجية,'' Zenodo (2026).
		\bibitem{BCLMEC} A.\,V.\ Yakushev, ``ملحق BCLM‑EC: الساعات الجينية,'' Zenodo (2026).
	\end{thebibliography}
	
\end{document}